\section{Submersions, Immersions, and Embeddings}

\subsection{Maps of constant rank}

\bdefn

    Let $F\colon M\to N$ be smooth, and $p\in M$.
    Then $F$'s {\emphcolor rank} at $p$ is the rank of the linear transformation $dF_p\colon T_pM\to T_{Fp}N$.

\edefn

Taking the coordinate bases of the tangent spaces, $F$'s rank is just the rank of the Jacobian of $F$ at $p$.
Equivalently it is the dimension of $dF_p$'s image in $T_{Fp}N$.
A matrices rank is bound by the dimensions of its domain and codomain, so $\rank_p F\leq\min\set{\dim M,\dim N}$.
If $F$'s rank at $p$ is equal to this minimum, we say that $F$ has {\emphcolor full rank} at $p$.

\bdefn

    $F\colon M\to N$ has {\emphcolor constant rank} if its rank is equal at every point in $M$.
    $F$ is a {\emphcolor smooth submersion} if its differential is surjective at every point, equivalently if
    $F$ has constant rank $\dim N$.
    $F$ is a {\emphcolor smooth immersion} if its differential is injective at every point, equivalently if
    $F$ has constant rank $\dim M$.

\edefn

\bprop

    Let $F\colon M\to N$ be smooth and at $p\in M$, $F$ has full rank.
    Then there is a neighborhood $U$ of $p$ such that $F|_U$ has full constant rank.
    That is to say, if $F$ has full rank at a point, there is a neighborhood where $F$ is a smooth immersion or
    submersion.

\eprop

\bproof

    Let $(U,\phi)$ and $(V,\psi)$ be charts of $M$ and $N$ respectively, and let $\hat F$ be $F$'s coordinate representation.
    Then the map $p\mapsto J\hat F_{\phi(p)}$ is continuous.
    The space of full-rank matrices is open, so its preimage under this map is open.
    This gives a neighborhood of $p$ where $F$ has full rank.
    \qed

\eproof

\bexam

    Let $M=M_1\times\cdots\times M_n$ be the direct product of smooth manifolds, and consider the projection
    $\pi=\pi_i\colon M\to M_i$: this is a smooth submersion.
    Indeed, for $p\in M$ let $(U_j,\phi_j)$ be charts containing $p^i$, so that $\pi$'s representation is
    $$ \hat\pi = \phi_i\circ\pi\circ(\phi_1^{-1},\dots,\phi_n^{-1}) , $$
    which maps $\bar x$ to $x^i$: i.e.\ $\hat\pi$ is the projection of $\bR^{n_1+\cdots+n_k}\to\bR^{n_i}$.
    The Jacobian of this has full rank, as desired.

\eexam

\bexam

    Let $M$ be smooth, and $\pi\colon TM\to M$ is a smooth submersion.
    Indeed, if $(U,\phi)$ is a chart of $M$ and $(\pi^{-1}(U),\tilde\phi)$ is its corresponding chart in $TM$, then
    $$ \hat\pi = \phi\circ\pi\circ\tilde\phi^{-1}\colon (\bar x,\bar v) \mapsto \bar x , $$
    which clearly has full rank.

\eexam

\subsection{Local diffeomorphisms}

\bdefn

    A smooth map $F\colon M\to N$ is a {\emphcolor local diffeomorphism} if for every $p\in M$, there is a neighborhood
    $p\in U$ such that $F(U)\subseteq N$ is open and $F|_U\colon U\to F(U)$ is a diffeomorphism.

\edefn

\bthrm

    Suppose $F\colon M\to N$ is smooth between manifolds without boundary.
    If $dF_p$ is invertible at some $p\in M$, there are connected neighborhoods $p\in U_0$ and $F(p)\in V_0$ such that
    $F|_{U_0}\colon U_0\to V_0$ is a diffeomorphism.

\ethrm

\bproof

    $dF_p$ being an isomorphism implies $M$ and $N$ have the same dimension.
    Take a chart centered at $p$, $(U,\phi)$, and a chart containing $F(U)$, $(V,\psi)$, centered at $F(p)$.
    Then $\hat F=\psi\circ F\circ\phi^{-1}\colon\phi(U)\to\psi(V)$ has $\hat F(\hat p)=0$ and $\hat F$'s Jacobian
    at $\hat p$ is invertible.
    The inverse function theorem (the ordinary Euclidean one) implies that there are connected neighborhoods
    $\hat U_0\subseteq\phi(U)$ and $\hat V_0\subseteq\psi(V)$ where $\hat F|_{\hat U_0}\colon\hat U_0\to\hat V_0$ is
    a diffeomorphism.
    Lifting these to $M$ and $N$ shows that $F|_{U_0}\colon U_0\to V_0$ is a diffeomorphism for $U_0=\phi^{-1}\hat U_0$
    and $V_0=\psi^{-1}\hat V_0$.
    \qed

\eproof

\bcoro

    A smooth map between manifolds without boundary is a local diffeomorphism iff its differential is invertible at every
    point.

\ecoro

The following are obvious properties of local diffeomorphisms.

\bprop

    \benum
        \item The composition of local diffeomorphisms is a local diffeomorphism;
        \item Every finite product of local diffeomorphisms is a local diffeomorphism;
        \item Every local diffeomorphism is a local homeomorphism and an open map;
        \item The restriction of a local diffeomorphism to an open submanifold with or without boundary is a local
        diffeomorphism;
        \item Every diffeomorphism is a local diffeomorphism;
        \item A bijective local diffeomorphism is a diffeomorphism;
    \eenum

\eprop

\bproof

    The only non-obvious claim is the last.
    Suppose $F\colon M\to N$ is a bijective local diffeomorphism, then $F^{-1}\colon N\to M$ is smooth.
    This is because for every $F(p)\in N$, there is a neighborhood $F(p)\in V$ where $F^{-1}|_V$ is smooth.
    Thus by the gluing lemma, $F^{-1}$ is smooth.
    \qed

\eproof

\bprop

    Let $F\colon M\to N$ be a smooth map between manifolds without boundary.
    \benum
        \item $F$ is a local diffeomorphism iff it is a smooth immersion and a smooth submersion.
        \item If $\dim M=\dim N$ then if $F$ is either a smooth immersion or submersion, then $F$ is a local diffeomorphism.
    \eenum

\eprop

\bproof

    The first point is clear, and if $F$ is a smooth immersion or submersion, then it must be the other (as injectivity
    or surjectivity of a linear map between spaces of the same dimension must be invertible).
    \qed

\eproof

\bexam

    Recall the map $\epsilon\colon\bR\to S^1$ by $\epsilon(t)=\exp(2\pi it)$.
    Then $\epsilon$'s representations are $t\mapsto\cos(2\pi t),\sin(2\pi t)$ which are local diffeomorphisms.
    So $\epsilon$ is a local diffeomorphism, and its $n$-fold product $\epsilon^n\colon\bR^n\to T^n$ is as well.

\eexam

\bthrm[title=Rank theorem]

    Suppose $M,N$ are smooth manifolds without boundary with dimensions $m,n$, and $F\colon M\to N$ is a map of constant
    rank $r$.
    Then for every $p\in M$ there exist smooth charts $(U,\phi)$ centered at $p$ and $(V,\psi)$ centered at $F(p)$
    containing $F(U)$, relative to which $F$ has coordinate representation
    $$ \hat F(x^1,\dots,x^r,x^{r+1},\dots,x^n) = (x^1,\dots,x^j,0,\dots,0) . $$
    In particular, if $F$ is a smooth submersion then $\hat F\colon\bR^n\to\bR^r$ is projection, and if $F$ is a smooth
    immersion then $\hat F\colon\bR^n\to\bR^r$ is inclusion (appending zeroes).

\ethrm

\bproof

    Without loss of generality, we can assume $F\colon U\to V$ is a smooth map between open subsets of Euclidean space,
    $U\subseteq\bR^n,V\subseteq\bR^m$.
    The Jacobian $JF(0)$ having rank $r$ means that it has some $r\times r$ minor with nonzero determinant.
    By reordering coordinates we can assume this is the upper-left minor.
    Let us label the standard coordinates as $(x,y)\in\bR^r\times\bR^{m-r}$ in the domain and $(v,w)\in\bR^r\times\bR^{n-r}$
    in the codomain.

    Then $F(x,y)=(Q(x,y),R(x,y))$ for smooth maps $Q\colon U\to\bR^r$ and $R\colon U\to\bR^{n-r}$, and our assumption
    means that $JQ(0)$ is non-singular.
    Define $\phi\colon U\to\bR^m$ by $\phi(x,y)=(Q(x,y),y)$ so that its Jacobian is
    $$ J\phi(0,0) = \pmatrix{\pdvof{Q^i}{x^j}(0) & \pdvof{Q^i}{y^j}(0) \cr 0 & \delta_{ij}} . $$
    This is non-singular and thus $\phi$ restricts to a diffeomorphism in a neighborhood of $0$, say $U_0$ and
    $\phi(U_0)=\tilde U_0$.
    We can assume that $\tilde U_0$ is an open cube, potentially by shrinking these two sets.

    Write $\phi^{-1}(x,y)=(A(x,y),B(x,y))$ for $A\colon\tilde U_0\to\bR^r$ and $B\colon\tilde U_0\to\bR^{m-r}$.
    Then we have
    $$ (x,y) = \phi(A(x,y),B(x,y)) = (Q(A(x,y),B(x,y)),B(x,y)) , $$
    so $B(x,y)=y$ and $\phi^{-1}$ has the form $(x,y)\mapsto(A(x,y),y)$.
    On the other hand, we have
    $$ (x,y) = \phi^{-1}\phi(x,y) = \phi^{-1}(Q(x,y),y) = (A(Q(x,y),y),y) , $$
    and so $A(Q(x,y),y)=x$ and $Q(A(x,y),y)=x$.

    Thus
    $$ F\circ\phi^{-1}(x,y) = (Q(A(x,y),y),R(A(x,y),y)) = (x,\tilde R(x,y)) $$
    for $\tilde R(x,y)=R(A(x,y),y)$.
    The Jacobian of this map is
    $$ J(F\circ\phi^{-1}) = \pmatrix{\delta_{ij} & 0\cr \pdvof{\tilde R^i}{x^j} & \pdvof{\tilde R^i}{y^j}} , $$
    which has rank $r$ as the product of $F$ and $\phi^{-1}$'s Jacobians, the latter of which is invertible.
    The first $r$ columns here are obviously linearly independent, and thus the rank of this map is $r$ iff
    $\ipdvof{\tilde R^i}{y^j}$ are zero on $\tilde U_0$.
    Since $\tilde U_0$ is a cube, this implies $\tilde R$ is independent of $(y^1,\dots,y^{m-r})$ so that if we define
    $S(x)=\tilde R(x,0)$ then $R(x,y)=S(x)$ for all $y$ and we have
    $$ F\circ\phi^{-1}(x,y) = (x,S(x)) . $$

    Now we want to define a smooth chart in some neighborhood of $0$ in $V$ such that
    $$ \psi\circ F\circ\phi^{-1} = \psi(x,S(x)) = (x,0) . $$
    Let $V_0=\set{(v,w)\in V}[(v,0)\in\tilde U_0]$, which is open.
    Then $F\circ\phi^{-1}(\tilde U_0)\subseteq V_0$ since $\tilde U_0$ is an open cube, and thus $F(U_0)\subseteq V_0$.
    Define $\psi\colon V_0\to\bR^n$ by $\psi(v,w)=(v,w-S(v))$, which has inverse $\psi^{-1}(x,y)=(s,t+S(s))$.
    And then
    $$ \psi\circ F\circ\phi^{-1}(x,y) = \psi(x,S(x)) = (x,0) $$
    as desired.
    \qed

\eproof

\bcoro

    Let $F\colon M\to N$ be a smooth map between manifolds without boundary, and $M$ is connected.
    Then $F$ has constant rank iff for every $p\in M$, $F$ has a coordinate representation which is linear.

\ecoro

\bproof

    If $F$ has a linear representation in a neighborhood of each point, then $F$ has constant rank in a neighborhood
    of each point.
    By connectedness, $F$ has constant rank.

    Conversely, if $F$ has constant rank, then apply the rank theorem as above.
    \qed

\eproof

\bthrm[title=Global rank theorem]

    If $F\colon M\to N$ is as above of constant rank, then
    \benum
        \item if $F$ is surjective, it is a smooth submersion;
        \item if $F$ is injective, it is a smooth immersion;
        \item if $F$ is bijective, it is a diffeomorphism.
    \eenum

\ethrm

\bproof

    Let $m,n$ be the dimensions of $M,N$ respectively, and $r$ the rank of $F$.
    If $F$ is not a smooth submersion then $r<n$.
    Consider then a coordinate representation of $F$ with respect to $(U,\phi),(V,\psi)$ which is a projection.
    Without loss of generality, $F(\cl U)\subseteq V$.

    Then $F(\cl U)$ is a compact subset of $\set{y\in V}[y^{r+1}=\cdots=r^n=0]$, and is thus closed in $N$.
    But it can't contain any open subset of $N$, and thus is nowhere dense.
    If we choose countably many charts covering $M$, $(U_i,\phi_i)$, then $F(M)$ is the union of the nowhere dense sets
    $F(\cl U_i)$, and thus has empty interior.
    So $F$ cannot be surjective.

    If $F$ is not a smooth immersion, then for $p\in M$ we have a coordinate representation as in the rank theorem.
    But such a coordinate representation is not injective, so $F$ cannot be.

    The last point is because the previous two points imply $F$ is a bijective local diffeomorphism, thus it is a
    diffeomorphism.
    \qed

\eproof

\subsection{Embeddings}

\bdefn

    A {\emphcolor smooth embedding} is a smooth map $F\colon M\to N$ which is a topological embedding (a homeomorphism onto
    its image), and a smooth immersion.

\edefn

A smooth embedding is not just a smooth topological embedding!

\bexam

    The inclusion of an open subset of a manifold in the manifold is a smooth embedding.

\eexam

\bexam

    For $M=M_1\times\cdots\times M_n$ and $p_i\in M_i$, the maps
    $$ \iota_j \colon M_j\to M,\qquad \iota_j(q) = (p_1,\dots,q,\dots,p_n) $$
    are smooth embeddigs.
    They can be given coordinate representations $\bar x\mapsto(\bar x,0)$ which is clearly a smooth embedding.

\eexam

\bprop

    Suppose $M,N$ are smooth manifolds and $F\colon M\to N$ is an injective smooth immersion.
    If any of the following hold, it is a smooth embedding.
    \benum
        \item $F$ is open or closed;
        \item $F$ is proper (preimages of compact subsets are compact);
        \item $M$ is compact;
        \item $M$ has empty boundary and $\dim M=\dim N$.
    \eenum

\eprop

\bproof

    If $F$ is open or closed and injective, it is a topological embedding, proving (1).
    If $F$ is proper, it is closed.
    If $M$ is compact, then $F$ is proper (since preimages of compact subspaces are closed in $M$ and thus compact).

    Now suppose $M$ has empty boundary and $\dim M=\dim N$.
    Since $dF_p$ is an isomorphism, we must have $F(M)\subseteq\Int N$.
    Thus $F\colon M\to\Int N$ is a local diffeomorphism, and thus open.
    So $F$ is a composition of open maps, and by (1) is a smooth embedding.
    \qed

\eproof

\bthrm[title=Local embedding theorem]

    If $M,N$ are manifolfds with or without boundary and $F\colon M\to N$ is smooth, then $F$ is a smooth immersion
    iff every point in $M$ has a neighborhood $U$ such that $F|_U\colon U\to N$ is a smooth embedding.

\ethrm

\bproof

    If every point has a neighborhood where $F$ is a smooth embedding, then $F$ has full rank and is thus a smooth
    immersion.

    Conversely if $F$ is a smooth immersion, then the rank theorem shows that $F$ has a coordinate representation
    of the form $F(\bar x)=(\bar x,0)$, which is injective, so $F|_U$ is injective.
    So we have that for some neighborhood $F|_U$ is injective.
    Now if $V\subseteq U$ is precompact, then $F|_{\cl V}$ is a topological embedding as it is injective and proper.
    Thus $F|_V$ is a topological embedding and a smooth immersion, as desired.
    \qed

\eproof

\subsection{Embedded and immersed submanifolds}

\bdefn

    Let $M$ be a manifold, then an {\emphcolor embedded submanifold} is a subset $S\subseteq M$ which is a manifold without boundary
    such that the inclusion $S\inj M$ is a smooth embedding.
    For an embedded submanifold $S\subseteq M$, its {\emphcolor codimension} is $\dim M-\dim S$.

\edefn

\bnote

    Embedded submanifolds are also called {\emphcolor regular submanifolds}.

\enote

Note that $S\inj M$ being a smooth embedding is the same as saying that it is a smooth immersion, and $S$ has the subspace topology.

\bprop

    The embedded submanifolds of $M$ with codimension $0$ are precisely the open subsets of $M$.

\eprop

\bproof

    For $\iota\colon U\inj M$ the inclusion map, we need only show it is a smooth immersion.
    And $\iota$'s coordinate representation relative to any point in $U$ is just the identity (take the same chart), which is
    a smooth immersion.

    Conversely, if $\iota\colon U\inj M$ is a smooth embedding of codimension $0$, then it is a local diffeomorphism.
    Thus it is open, and so $U$ must be an open subset of $M$.
    \qed

\eproof

\bprop

    Let $F\colon M\to N$ be a smooth embedding, then $S=F(M)$ has a unique smooth structure making it an embedded submanifold of $N$
    diffeomorphic to $M$.

\eprop

\bproof

    We need to give $S$ the subspace topology relative to $N$, and then by assumption it is homeomorphic to $M$.
    We need to lift the smooth structure of $M$ to $S$, that is take charts $(F(U),\phi\circ F^{-1})$ for a chart $(U,\phi)$ for $M$.
    This is a smooth structure, and it turns $F$ into a diffeomorphism onto $S$, and $\iota\colon S\inj N$ into a smooth embedding.
    \qed

\eproof

\bexam

    Consider the product manifold $M\times N$ where $M$ doesn't have boundary.
    Then for $p\in N$, $M\times\set p$ is an embedded submanifold diffeomorphic to $M$.
    Indeed, $F\colon M\to M\times N$ given by $x\mapsto(x,p)$ is a smooth embedding with image $M\times\set p$.

    This embedded submanifold is called a {\emphcolor slice} of the product manifold.

\eexam

\bdefn

    Let $M$ be a smooth manifold, then an {\emphcolor immersed submanifold} is a subset $S\subseteq M$ endowed with a
    smooth structure making the inclusion $\iota\colon S\inj M$ a smooth immersion.
    The {\emphcolor codimension} of an immersed submanifold is defined still to be $\dim M-\dim S$.

\edefn

Since the inclusion is not necessarily a topological embedding, $S$'s topology need not be the subspace topology.

\bprop

    Suppose $F\colon M\to N$ is an injective smooth immersion, and let $S=F(M)$.
    Then $S$ has a unique topology and smooth structure making it an immersed submanifold of $N$, and $F\colon M\to S$
    a diffeomorphism onto its image.

\eprop

\bproof

    The proof is similar to the one above.
    Give $S$ the topology where $U\subseteq S$ is open iff $F^{-1}U\subseteq M$ is.
    Then charts $(F(U),\phi\circ F^{-1})$ for a chart $(U,\phi)$ for $M$.
    \qed

\eproof

\bprop

    Suppose $S\subseteq M$ is an immersed submanifold.
    Then if any of the following hold, $S$ is embedded:
    \benum
        \item $S$ has codimension $0$,
        \item The inclusion map $S\inj M$ is proper (preimages of compact sets are compact),
        \item $S$ is compact.
    \eenum

\eprop

\bproof

    If $S$ is compact, then the inclusion map is proper (closed subsets of compact sets are compact), so (2) implies (3).
    We need only prove the first two points.

    If $S$ has codimension $0$, then a smooth immersion is a smooth submersion is a local diffeomorphism.
    Thus the inclusion map is an injective local diffeomorphism, thus a topological embedding.
    \qed

\eproof

\bprop

    If $S\subseteq M$ is an immersed submanifold, then for every $p\in S$ there is a neighborhood $p\in U\subseteq S$
    which is an embedded submanifold of $M$.

\eprop

\bproof

    We showed above that smooth immersions are locally smooth embeddings.
    \qed

\eproof

Note that if $S\subseteq M$ is an immersed submanifold and $F\colon M\to N$ is smooth, then
$$ F\circ\iota = F|_S\colon S\to N $$
is smooth.
Thus we can restrict smooth maps to immersed submanifolds.
But we can't always restrict the codomains of smooth maps, as the topology of an immersed submanifold differs, and thus the restriction
may not even be continuous.

\bthrm

    Suppose $F\colon M\to N$ is smooth and $S\subseteq N$ is an immersed submanifold containing $F$'s image.
    If the restriction $F\colon M\to S$ is continuous, it is smooth.

\ethrm

\bproof

    Let $p\in M$ arbitrary, and let $q=f(p)\in S$.
    Then $\iota\colon S\to N$ locally has the representation $\bar x\mapsto(\bar x,0)$, let $(V,\psi)\subseteq S$ and $(W,\tilde\psi)\subseteq N$
    be charts centered at $q$ giving this representation.
    This means that $\tilde\psi\circ\iota\circ\psi^{-1}\colon\bar x\mapsto(\bar x,0)$.
    If we let $\pi$ be the projection $(\bar x,\bar y)\mapsto\bar x$, then we have
    $$ \pi\circ\tilde\psi\circ\iota\circ\psi^{-1} = \id \implies \psi = \pi\circ\tilde\psi\circ\iota . $$

    Since $F$ is continuous, $U=F^{-1}(V_0)$ is an open neighborhood of $p$ in $M$.
    Choose a smooth chart $(U_0,\phi)$ centered at $p$ with $p\in U_0\subseteq U$.
    Then if we denote by $\tilde F\colon M\to S$ the restriction onto the codomain, we have $\iota\circ\tilde F=F$, and
    $$ \psi\circ\tilde F\circ\phi^{-1} = \pi\circ\tilde\psi\circ\iota\circ\tilde F\circ\phi^{-1} = \pi\circ\tilde\psi\circ F\circ\phi^{-1} , $$
    which is smooth, as $F$ and $\pi$ are.
    \qed

\eproof

\bcoro

    Suppose $F\colon M\to N$ is smooth and $S\subseteq N$ is an embedded submanifold containing $F$'s image.
    Then the restriction of $F$ to $M\to S$ is smooth.

\ecoro

\bproof

    The restriction of a continuous map to its image is continuous with respect to the subspace topology, which $S$ has.
    \qed

\eproof

Note that if $S\subseteq M$ is an immersed submanifold, then by definition $d\iota_p\colon T_pS\to T_pM$ is an injection for every $p\in S$.
Thus we can and will identify $T_pS$ with its image in $T_pM$.
Recall that for $f\in C^\infty(M)$ and $v\in T_pS$,
$$ d\iota_pv(f) = v(f\circ\iota) = v(f|_S) . $$

\bprop

    Let $S\subseteq M$ be an immersed submanifold, with a point $p\in S$.
    Then the tangent space $T_pS\subseteq T_pM$ is characterized by
    $$ T_pS = \set{v\in T_pM}[\hbox{$vf=0$ for all smooth real-valued functions $f$ which are zero on $S$}] . $$

\eprop

\bproof

    If $v\in T_pS$ then clearly $d\iota_pv(f)=v(f|_S)=0$ if $f|_S=0$.
    So clearly the lefthand side is contained in the right.
    Conversely if $v(f)=0$ for all $f|_S=0$, then we need to find a origin for $v$ in $T_pS$.

    Let $(U,\phi)\subseteq S$ and $(V,\psi)\subseteq M$ be charts centered at $p$ for which $\iota$'s coordinate representation is
    $\bar x\mapsto(\bar x,0)$ for $\bar x=(x^1,\dots,x^k)$.
    Then we see that $T_pS\subseteq T_pM$ is spanned by $\ievpdv{x^1}_p,\dots,\ievpdv{x^k}_p$.
    Now we write $v$ in terms of the entire basis set,
    $$ v = \sum_{i=1}^nv^i\evpdv{x^i}_p . $$
    Consider a bump function $h$ supported in $U$ which is equal to $1$ in a neighborhood of $p$, and define $f(x)=h(x)x^j$ for $j>k$.
    Then $f$ vanishes on $S$, but is identically $1$ in a neighborhood of $p$ so that
    $$ 0 = v(f) = \sum_{i=1}^n v^i\pdvof{h(x)x^j}{x^i}(p) = v^j , $$
    as desired.
    \qed

\eproof

\bdefn

    Let $\Phi\colon M\to N$ be any map and $c\in N$.
    Then $\Phi$'s preimage on $c$, $\Phi^{-1}c$, is called a {\emphcolor level set} of $\Phi$.

\edefn

\bprop

    Let $\Phi\colon M\to N$ be smooth with constant rank $r$.
    Then each level set of $\Phi$ is an embedded submanifold of $M$ of codimension $r$.

\eprop

\bproof

    Let $m=\dim M,n=\dim N,k=m-r$, and choose $c\in N$.
    Let $S=\Phi^{-1}c\subseteq M$ be the level set at $c$.
    Then for every $p\in S$, there are smooth charts centered at $p\in (U,\phi)$ and $c\in(V,\psi)$ such that $\Phi$ has the representatiion
    $$ \psi\circ\Phi\circ\phi^{-1}\colon (x^1,\dots,x^m) \mapsto (x^1,\dots,x^r,0,\dots,0) . $$
    Then for $q\in U\cap S$, we have that $\psi\circ\Phi(q)=\psi(c)=0$ meaning $\phi(q)=(x^1,\dots,x^r,x^{r+1},\dots,x^m)$ where $x^1,\dots,x^r=0$.
    Thus we can define $\tilde\phi\colon U\cap S\to\bR^k$ as the projection of $\phi$ onto its last $k$ coordinates.
    These are homeomorphisms as they just map $\phi(S\cap U)\cong 0\times\bR^k\cong\bR^k$, and they are smoothly compatible.
    Thus they form a smooth atlas for $S$.

    The inclusion is a smooth embedding, as $\phi\circ\iota\circ\tilde\phi^{-1}\colon\bR^k\to\bR^n$ is just the inclusion and thus has full rank.
    \qed

\eproof

\bcoro

    If $\Phi\colon M\to N$ is a smooth submersion, then each level set of $\Phi$ is an embedded submanifold of $M$ with codimension $\dim N$.

\ecoro

\bcoro

    If $\Phi\colon M\to N$ is smooth, then a point $p\in M$ is a {\emphcolor regular point} of $\Phi$ if $\Phi$'s rank at $p$ is $\dim N$,
    and otherwise is a {\emphcolor critical point}.

    A point $c\in N$ is a {\emphcolor regular value} if each $p\in\Phi^{-1}c$ is a regular value, and a {\emphcolor critical value} otherwise.

    If $c\in N$ is a regular value, then $\Phi^{-1}c$ is called a {\emphcolor regular level set}.
    That is, it is a level set all of whose points are regular.

\ecoro

Note that if $\Phi^{-1}c=\emptyset$, then $c$ is a regular value.
If $\Phi$ is a smooth submersion, then every point is regular (these are of course equivalent), and thus all of its level sets are regular.
We can then improve the above corollary.

\bcoro

    Every regular level set of a smooth map is an embedded submanifold of the domain whose codimension is equal to the dimension of the codomain.

\ecoro

\bproof

    Let $\Phi\colon M\to N$ be smooth, and $c\in N$ a regular value.
    Let $U\subseteq M$ be the set of points $p\in M$ where $d\Phi_p=\dim N$.
    For $p\in U$ we can take a coordinate representation of $\Phi$ in a neighborhood of $U$, call it $\hat\Phi$.
    Then $J\hat\Phi(p)$ has full rank, and then must have full rank in a neighborhood of $p$, since $q\mapsto J\hat\Phi(q)$ is continuous
    and the set of matrices of full rank is open.
    Thus $U$ is an open set, and thus an embedded submanifold of $M$, and contains $\Phi^{-1}c$ by assumption.

    Thus $\Phi|_U\colon U\to N$ is a smooth submersion, and by the above corollary $\Phi^{-1}c\subseteq U$ is an embedded submanifold.
    An embedded submanifold of an embedded submanifold is an embedded submanifold (the composition of smooth embeddings is a smooth embedding), so $\Phi^{-1}c$
    is an embedded submanifold of $M$.
    \qed

\eproof

\subsection{Submersions}

A {\emphcolor section} of a map $\pi\colon M\to N$ is a map $\sigma\colon N\to M$ which is a right inverse: $\pi\circ\sigma=\id_N$.
If our maps are continuous or smooth, we speak of continuous or smooth sections.
A {\emphcolor local section} is a continuous map $\sigma\colon U\to M$ defined in an open subset of $N$ where $\pi\circ\sigma_U=\id_U$.

\bthrm

    If $\pi\colon M\to N$ is smooth, then it is a smooth submersion iff every point of $M$ is in the image of a smooth local section.

\ethrm

\bproof

    If $\pi$ is a smooth submersion, then it has the local representation $(x^1,\dots,x^m)\mapsto(x^1,\dots,x^n)$ around any point $p\in M$.
    For $\epsilon>0$ small enough, the coordinate cube $C_\epsilon=\set{x}[\abs{x^i}<\epsilon\hbox{ for $i=1,\dots,m$}]$ is an open neighborhood
    centered at $p$ whose image under $\pi$ is the coordinate cube $C'_\epsilon=\set{y}[\abs{y^i}<\epsilon\hbox{ for $i=1,\dots,n$}]$.
    Then defining $\sigma\colon C'_\epsilon\to C_\epsilon$ by $\sigma(y)=(y,0)$ gives a local section of $\pi$ and $p=\sigma(0)$ as desired.

    Conversely let $\sigma\colon U\to M$ be a local section with $p$ in its image, say $p=\sigma(q)$.
    Then since $\pi\circ\sigma=\id_U$, we have that $d\pi_p\circ d\sigma_q=\id_{T_qN}$, which means $d\pi_p$ must be surjective.
    \qed

\eproof

This allows us to extend the definition of a smooth submersion to a topological one: a {\emphcolor topological submersion} is a continuous map
$X\to Y$ such that every point of $X$ is in the image of a local section of the map.
This theorem shows that a smooth submersion is just a smooth topological submersion.

\bdefn

    Recall that a continuous map $\pi\colon X\to Y$ is a {\emphcolor quotient map} if it is surjective and $Y$'s topology is the quotient topology
    induced by $\pi$.
    That is, $U\subseteq Y$ is open iff $\pi^{-1}U\subseteq X$ is.

\edefn

Note that a quotient map is just a surjective, continuous, and open (equivalently closed) map.

\bprop

    Let $\pi\colon M\to N$ be a smooth submersion.
    Then $\pi$ is an open map, and if it is surjective it is a quotient map.

\eprop

\bproof

    Let $W\subseteq M$ open, and $q\in\pi(W)$.
    Then for any $p\in W$ such that $\pi(p)=q$, there is a neighborhood $U$ of $q$ and a local section $\sigma\colon U\to N$ such that $\sigma(q)=p$.
    (Note that $p$ being in the image of $\sigma$ requires $\sigma(q)=p$, since if $\sigma(r)=p$ then $r=\pi\sigma(r)=\pi(p)=q$.)
    Then for $y\in\sigma^{-1}(W)$, since $\sigma(y)\in W$ we have that $y\in\pi(W)$.
    And then $\sigma^{-1}(W)$ is an open subset of $\pi(W)$ containing $q$, and so $\pi(W)$ is open.

    Thus $\pi$ is an open map.
    Surjective open continuous maps are quotient maps.
    \qed

\eproof

\bthrm

    Let $\pi\colon M\to N$ be a surjective smooth submersion.
    Then for any manifold $P$, a map $F\colon P\to N$ is smooth iff $F\circ\pi\colon M\to P$ is smooth.

\ethrm

\bproof

    If $F$ is smooth then so is $F\circ\pi$ is.
    Conversely if $F\circ\pi$ is smooth, let $q\in N$ and since $\pi$ is surjective, there is a $p\in M$ with $\pi(p)=q$.
    And then there is a local section $\sigma\colon U\to M$ with $U$ a neighborhood of $q$, and $\sigma(q)=p$.
    Then $\pi\circ\sigma=\id_U$ implies
    $$ F|_U = F|_U\circ\id_U = F|_U\circ(\pi\circ\sigma) = (F\circ\pi)\circ\sigma, $$
    which is smooth.
    So for every point in $N$, there is a neighborhood around it in which $F$ is smooth.
    Thus $F$ is smooth.
    \qed

\eproof

\bthrm

    Let $\pi\colon M\to N$ be a surjective smooth submersion.
    If $F\colon M\to P$ is smooth which is constant on the fibers of $\pi$, there is a unique map $\tilde F\colon N\to P$ such that
    $F=\tilde F\circ\pi$.

\ethrm

\bproof

    As a quotient space, there is a continuous $\tilde F$ satisfying this.
    Then since $\tilde F\circ\pi$ is smooth, so is $\tilde F$.
    \qed

\eproof

